\subsection{Matrix Addition \& Multiplication}

An $m \times n$ matrix $\mathbf{A}$ has entries $a_{ij}$ with $i$ the row index and $j$ the column index:
\[
  \mathbf{A}
  = \begin{bmatrix}
    a_{11} & a_{12} & \cdots & a_{1n} \\
    a_{21} & a_{22} & \cdots & a_{2n} \\
    \vdots & \vdots & \ddots & \vdots \\
    a_{m1} & a_{m2} & \cdots & a_{mn}
  \end{bmatrix}
  = [a_{ij}].
\]

\subsubsection{Trivial Addition and Scalar Multiplication}

\begin{definition}{Matrix Addition}{mat-add}
Let $\mathbf{A}, \mathbf{B} \in \R^{m \times n}$. Their \textbf{sum} $\mathbf{A} + \mathbf{B} \in \R^{m \times n}$ is defined entrywise by
\[
  (\mathbf{A} + \mathbf{B})_{ij} = \mathbf{A}_{ij} + \mathbf{B}_{ij}.
\]
\end{definition}

\begin{definition}{Scalar Multiplication}{scalar-mul}
Let $\mathbf{A} \in \R^{m \times n}$ and $\lambda \in \R$. Then $\lambda \mathbf{A} \in \R^{m \times n}$ is defined entrywise by
\[
  (\lambda \mathbf{A})_{ij} = \lambda \mathbf{A}_{ij}.
\]
\end{definition}

\textbf{Properties.} For $\mathbf{A}, \mathbf{B}, \mathbf{C} \in \R^{m \times n}$ and scalars $\lambda, \delta \in \R$:
\begin{itemize}
  \item Commutative: $\mathbf{A} + \mathbf{B} = \mathbf{B} + \mathbf{A}$.
  \item Associative: $(\mathbf{A} + \mathbf{B}) + \mathbf{C} = \mathbf{A} + (\mathbf{B} + \mathbf{C})$.
  \item Distributive: $\lambda(\mathbf{A} + \mathbf{B}) = \lambda \mathbf{A} + \lambda \mathbf{B}$ and $(\lambda + \delta)\mathbf{A} = \lambda \mathbf{A} + \delta \mathbf{A}$.
  \item Additive identity: $\mathbf{A} + \mathbf{0} = \mathbf{A}$, where $\mathbf{0}$ is the $m \times n$ zero matrix.
\end{itemize}

\subsubsection{Matrix Multiplication}

\begin{definition}{Matrix Multiplication}{mat-mul}
Let $\mathbf{A} \in \R^{m \times n}$ and $\mathbf{B} \in \R^{n \times p}$ with $\mathbf{B}$ written by columns as $\mathbf{B} = \bigl[\,\mathbf{b}_1 \;\bigm|\; \mathbf{b}_2 \;\bigm|\; \cdots \;\bigm|\; \mathbf{b}_p\,\bigr]$. The \textbf{product} $\mathbf{A}\mathbf{B} \in \R^{m \times p}$ is
\[
  \mathbf{A}\mathbf{B}
  = \bigl[\,\mathbf{A}\mathbf{b}_1 \;\bigm|\; \mathbf{A}\mathbf{b}_2 \;\bigm|\; \cdots \;\bigm|\; \mathbf{A}\mathbf{b}_p\,\bigr].
\]
\end{definition}

\begin{example}
  Let
  $\mathbf{A} = \begin{bmatrix} 2 & 1 & -1 \\ 1 & 0 & -3 \end{bmatrix}$ ($2 \times 3$),
  $\mathbf{B} = \begin{bmatrix} 1 & -2 \\ 1 & 3 \\ 1 & 0 \end{bmatrix}$ ($3 \times 2$).
  Then $\mathbf{A}\mathbf{B}$ is $2 \times 2$:
  \[
    \mathbf{A}\mathbf{B}
    = \left[\,\mathbf{A}\begin{bmatrix} 1 \\ 1 \\ 1 \end{bmatrix} \;\bigm|\; \mathbf{A}\begin{bmatrix} -2 \\ 3 \\ 0 \end{bmatrix}\,\right]
    = \begin{bmatrix} 2 & -1 \\ -2 & -2 \end{bmatrix}.
  \]
\end{example}

\begin{remark}
Matrix multiplication is \textbf{not commutative}. Using the same $\mathbf{A}, \mathbf{B}$ as above, $\mathbf{B}\mathbf{A}$ is $3 \times 3$:
\[
  \mathbf{B}\mathbf{A}
  = \begin{bmatrix} 1 & -2 \\ 1 & 3 \\ 1 & 0 \end{bmatrix}
    \begin{bmatrix} 2 & 1 & -1 \\ 1 & 0 & -3 \end{bmatrix}
  = \begin{bmatrix} 0 & 1 & 5 \\ 5 & 1 & -10 \\ 2 & 1 & -1 \end{bmatrix}.
\]
$\mathbf{A}\mathbf{B}$ and $\mathbf{B}\mathbf{A}$ are not even the same size.
\end{remark}

\subsubsection{Row-Column View of Matrix Multiplication}

The $(i, j)$ entry of $\mathbf{A}\mathbf{B}$ comes from row $i$ of $\mathbf{A}$ paired with column $j$ of $\mathbf{B}$, and lives at row $i$, column $j$ of $\mathbf{A}\mathbf{B}$:
\[
  \underbrace{\begin{bmatrix}
    i\; \text{row} \quad \vdots \\
    \boxed{a_{i1}\;\; a_{i2}\;\; \cdots\;\; a_{in}} \\
    \vdots
  \end{bmatrix}}_{\mathbf{A}}
  \;\times\;
  \underbrace{\begin{bmatrix}
    j\; \text{col} \quad \cdots &
    \boxed{\begin{matrix} b_{1j} \\ b_{2j} \\ \vdots \\ b_{nj} \end{matrix}} &
    \cdots
  \end{bmatrix}}_{\mathbf{B}}
  \;=\;
  \underbrace{\left[\begin{array}{c|c|c}
    & i\; \text{row} \quad \vdots & \\ \hline
    j\; \text{col} \quad \cdots & (\mathbf{A}\mathbf{B})_{ij} & \cdots \\ \hline
    & \vdots &
  \end{array}\right]}_{\mathbf{A}\mathbf{B}}
\]
Taking the dot product of these strips:
\[
  (\mathbf{A}\mathbf{B})_{ij}
  = \begin{bmatrix} a_{i1} & a_{i2} & \cdots & a_{in} \end{bmatrix}
    \begin{bmatrix} b_{1j} \\ b_{2j} \\ \vdots \\ b_{nj} \end{bmatrix}
  = a_{i1} b_{1j} + a_{i2} b_{2j} + \cdots + a_{in} b_{nj}.
\]

\begin{example}
  Compute $(\mathbf{B}\mathbf{A})_{33}$ for the $\mathbf{A}, \mathbf{B}$ above. Box row $3$ of $\mathbf{B}$ and column $3$ of $\mathbf{A}$:
  \[
    \underbrace{\begin{bmatrix}
      \begin{matrix} 1 & -2 \end{matrix} \\
      \begin{matrix} 1 & \phantom{-}3 \end{matrix} \\
      \boxed{\begin{matrix} 1 & \phantom{-}0 \end{matrix}}
    \end{bmatrix}}_{\mathbf{B}}
    \;\times\;
    \underbrace{\begin{bmatrix}
      \begin{matrix} 2 \\ 1 \end{matrix} &
      \begin{matrix} 1 \\ 0 \end{matrix} &
      \boxed{\begin{matrix} -1 \\ -3 \end{matrix}}
    \end{bmatrix}}_{\mathbf{A}}
    \;=\;
    \underbrace{\left[\begin{array}{cc|c}
      \cdot & \cdot & \vdots \\
      \cdot & \cdot & \vdots \\ \hline
      \cdots & \cdots & -1
    \end{array}\right]}_{\mathbf{B}\mathbf{A}}
  \]
  Taking the dot product:
  \[
    (\mathbf{B}\mathbf{A})_{33}
    = \begin{bmatrix} 1 & 0 \end{bmatrix}
      \begin{bmatrix} -1 \\ -3 \end{bmatrix}
    = (1)(-1) + (0)(-3) = -1.
  \]
\end{example}

\subsubsection{Properties of Matrix Multiplication}

For $\mathbf{A}$, $\mathbf{B}$, $\mathbf{C}$, and $\lambda \in \R$ (matrix sizes must be compatible):
\begin{itemize}
  \item Associative: $\mathbf{A}(\mathbf{B}\mathbf{C}) = (\mathbf{A}\mathbf{B})\mathbf{C}$.
  \item Not commutative in general (see counter-example above).
  \item Distributive: $(\mathbf{A} + \mathbf{B})\mathbf{C} = \mathbf{A}\mathbf{C} + \mathbf{B}\mathbf{C}$ and $\mathbf{C}(\mathbf{A} + \mathbf{B}) = \mathbf{C}\mathbf{A} + \mathbf{C}\mathbf{B}$.
  \item Scalar: $\lambda(\mathbf{A}\mathbf{B}) = (\lambda\mathbf{A})\mathbf{B} = \mathbf{A}(\lambda\mathbf{B})$.
  \item Identity: $\mathbf{A}\mathbf{I}_n = \mathbf{I}_m \mathbf{A} = \mathbf{A}$.
\end{itemize}

\subsection{Transpose}

\begin{definition}{Matrix Transpose}{transpose}
Let $\mathbf{A} \in \R^{m \times n}$. The \textbf{transpose} $\mathbf{A}^T \in \R^{n \times m}$ is the matrix whose rows are the columns of $\mathbf{A}$:
\[
  \mathbf{A} = \bigl[\,\mathbf{a}_1 \;\bigm|\; \mathbf{a}_2 \;\bigm|\; \cdots \;\bigm|\; \mathbf{a}_n\,\bigr]
  \quad\Longrightarrow\quad
  \mathbf{A}^T = \begin{bmatrix} \mathbf{a}_1 \\ \mathbf{a}_2 \\ \vdots \\ \mathbf{a}_n \end{bmatrix}.
\]
\end{definition}

\begin{example}
  \[
    \mathbf{A} = \begin{bmatrix} 1 & 2 & 3 \\ 0 & 3 & -5 \\ 1 & -1 & 2 \end{bmatrix}
    \quad\Longrightarrow\quad
    \mathbf{A}^T = \begin{bmatrix} 1 & 0 & 1 \\ 2 & 3 & -1 \\ 3 & -5 & 2 \end{bmatrix}.
  \]
  \[
    \mathbf{B} = \begin{bmatrix} 1 & 3 & -1 \\ 2 & 5 & 1 \end{bmatrix}
    \quad\Longrightarrow\quad
    \mathbf{B}^T = \begin{bmatrix} 1 & 2 \\ 3 & 5 \\ -1 & 1 \end{bmatrix}.
  \]
\end{example}

\textbf{Properties.} For matrices $\mathbf{A}, \mathbf{B}$ of compatible sizes:
\begin{itemize}
  \item $(\mathbf{A}^T)^T = \mathbf{A}$.
  \item $(\mathbf{A} + \mathbf{B})^T = \mathbf{A}^T + \mathbf{B}^T$.
  \item $(\mathbf{A}\mathbf{B})^T = \mathbf{B}^T \mathbf{A}^T$.
\end{itemize}

\begin{example}
  Let $\mathbf{x} = \begin{bmatrix} 3 \\ -2 \\ 2 \end{bmatrix}$. Then $\mathbf{x}^T = \begin{bmatrix} 3 & -2 & 2 \end{bmatrix}$, and
  \[
    \mathbf{x}^T \mathbf{x}
    = \begin{bmatrix} 3 & -2 & 2 \end{bmatrix} \begin{bmatrix} 3 \\ -2 \\ 2 \end{bmatrix}
    = \begin{bmatrix} 17 \end{bmatrix},
  \]
  \[
    \mathbf{x}\mathbf{x}^T
    = \begin{bmatrix} 3 \\ -2 \\ 2 \end{bmatrix} \begin{bmatrix} 3 & -2 & 2 \end{bmatrix}
    = \begin{bmatrix} 9 & -6 & 6 \\ -6 & 4 & -4 \\ 6 & -4 & 4 \end{bmatrix}.
  \]
\end{example}

\subsection{The Inverse Matrix}

\textbf{Motivation.} If $\mathbf{A}\mathbf{x} = \mathbf{b}$ and $\mathbf{A}$ has an inverse $\mathbf{A}^{-1}$, then solving a linear system is easy:
\[
  \mathbf{A}^{-1}\mathbf{A}\mathbf{x} = \mathbf{A}^{-1}\mathbf{b}
  \quad\Longrightarrow\quad
  \mathbf{x} = \mathbf{A}^{-1}\mathbf{b}.
\]

\begin{definition}{Matrix Inverse}{mat-inv}
Let $\mathbf{A} \in \R^{n \times n}$ be a square matrix. A matrix $\mathbf{B}$ is an \textbf{inverse} of $\mathbf{A}$ if
\[
  \mathbf{B}\mathbf{A} = \mathbf{I}_n
  \quad\text{and}\quad
  \mathbf{A}\mathbf{B} = \mathbf{I}_n.
\]
We denote $\mathbf{A}^{-1} = \mathbf{B}$.
\end{definition}

\begin{remark}
Some matrices don't have inverses, such as the zero matrix $\mathbf{0}$: for any $\mathbf{B}$, $\mathbf{0}\mathbf{B} = \mathbf{0} \neq \mathbf{I}_n$.
\end{remark}

\begin{example}
  Let $\mathbf{A} = \begin{bmatrix} 2 & 5 \\ 1 & 3 \end{bmatrix}$ and $\mathbf{B} = \begin{bmatrix} 3 & -5 \\ -1 & 2 \end{bmatrix}$. Then
  \[
    \mathbf{A}\mathbf{B}
    = \begin{bmatrix} 2 & 5 \\ 1 & 3 \end{bmatrix} \begin{bmatrix} 3 & -5 \\ -1 & 2 \end{bmatrix}
    = \begin{bmatrix} 1 & 0 \\ 0 & 1 \end{bmatrix},
    \quad
    \mathbf{B}\mathbf{A}
    = \begin{bmatrix} 3 & -5 \\ -1 & 2 \end{bmatrix} \begin{bmatrix} 2 & 5 \\ 1 & 3 \end{bmatrix}
    = \begin{bmatrix} 1 & 0 \\ 0 & 1 \end{bmatrix},
  \]
  so $\mathbf{B}$ is the inverse of $\mathbf{A}$.
\end{example}

\begin{theorem}{Uniqueness of Inverse}{inv-unique}
If $\mathbf{A}$ is invertible, then its inverse is unique.
\end{theorem}

\begin{proof}
By contradiction. Assume $\mathbf{B} \neq \mathbf{C}$ are both inverses of $\mathbf{A}$, so
\[
  \mathbf{B}\mathbf{A} = \mathbf{A}\mathbf{B} = \mathbf{I},
  \qquad
  \mathbf{C}\mathbf{A} = \mathbf{A}\mathbf{C} = \mathbf{I}.
\]
Then
\[
  \mathbf{C}
  = \mathbf{C}\mathbf{I}
  = \mathbf{C}(\mathbf{A}\mathbf{B})
  = (\mathbf{C}\mathbf{A})\mathbf{B}
  = \mathbf{I}\mathbf{B}
  = \mathbf{B},
\]
contradicting $\mathbf{B} \neq \mathbf{C}$.
\end{proof}

\begin{theorem}{$2 \times 2$ Inverse Formula}{inv-2x2}
Let $\mathbf{A} = \begin{bmatrix} a & b \\ c & d \end{bmatrix}$ with $ad - bc \neq 0$. Then
\[
  \mathbf{A}^{-1} = \frac{1}{ad - bc} \begin{bmatrix} d & -b \\ -c & a \end{bmatrix}.
\]
\end{theorem}

\begin{proof}
\[
  \mathbf{A}\mathbf{A}^{-1}
  = \begin{bmatrix} a & b \\ c & d \end{bmatrix} \cdot \frac{1}{ad - bc} \begin{bmatrix} d & -b \\ -c & a \end{bmatrix}
  = \frac{1}{ad - bc} \begin{bmatrix} ad - bc & 0 \\ 0 & ad - bc \end{bmatrix}
  = \mathbf{I}_2. \qedhere
\]
\end{proof}

\subsubsection{Elementary Matrices}

\begin{definition}{Elementary Matrix}{elem-mat}
An \textbf{elementary matrix} is a matrix obtained from an identity matrix by applying a single row operation.
\end{definition}

\begin{example}
  Starting from $\mathbf{I}_3 = \begin{bmatrix} 1 & 0 & 0 \\ 0 & 1 & 0 \\ 0 & 0 & 1 \end{bmatrix}$:
  \begin{align*}
    R_2 &\leftarrow 5 R_2: & \begin{bmatrix} 1 & 0 & 0 \\ 0 & 5 & 0 \\ 0 & 0 & 1 \end{bmatrix}, \\
    R_1 &\leftrightarrow R_3: & \begin{bmatrix} 0 & 0 & 1 \\ 0 & 1 & 0 \\ 1 & 0 & 0 \end{bmatrix}, \\
    R_3 &\leftarrow R_3 + b R_1: & \begin{bmatrix} 1 & 0 & 0 \\ 0 & 1 & 0 \\ b & 0 & 1 \end{bmatrix}.
  \end{align*}
\end{example}

\begin{theorem}{Row Operation as Elementary Matrix Multiplication}{row-op-elem}
Let $\mathbf{A} \in \R^{n \times n}$ and $\mathbf{E} \in \R^{n \times n}$ an elementary matrix. Then $\mathbf{E}\mathbf{A}$ is the matrix obtained from $\mathbf{A}$ by applying the same row operation that produced $\mathbf{E}$ from $\mathbf{I}_n$.
\end{theorem}

\begin{example}
  Swapping rows $1$ and $3$:
  \[
    \begin{bmatrix} 0 & 0 & 1 \\ 0 & 1 & 0 \\ 1 & 0 & 0 \end{bmatrix}
    \begin{bmatrix} 1 & 2 & 3 \\ 4 & 5 & 6 \\ 7 & 8 & 9 \end{bmatrix}
    = \begin{bmatrix} 7 & 8 & 9 \\ 4 & 5 & 6 \\ 1 & 2 & 3 \end{bmatrix}.
  \]
\end{example}

\subsubsection{Finding the Inverse via Row Reduction}

For $\mathbf{A} \in \R^{n \times n}$, if elementary matrices $\mathbf{E}_1, \mathbf{E}_2, \ldots, \mathbf{E}_k$ satisfy
\[
  \mathbf{E}_k \cdots \mathbf{E}_2 \mathbf{E}_1 \mathbf{A} = \mathbf{I}_n,
  \quad\text{then}\quad
  \mathbf{A}^{-1} = \mathbf{E}_k \cdots \mathbf{E}_2 \mathbf{E}_1.
\]
So finding $\mathbf{A}^{-1}$ becomes finding row operations that reduce $\mathbf{A}$ to $\mathbf{I}_n$.

\textbf{Procedure.} Form the augmented matrix $[\,\mathbf{A} \mid \mathbf{I}_n\,]$ and row reduce:
\[
  [\,\mathbf{A} \mid \mathbf{I}_n\,]
  \;\xrightarrow{\text{row ops}}\;
  [\,\mathbf{I}_n \mid \mathbf{A}^{-1}\,].
\]

\begin{example}
  Find the inverse of $\mathbf{A} = \begin{bmatrix} 2 & 1 & 1 \\ 0 & 1 & 2 \\ 0 & 0 & 2 \end{bmatrix}$.

  Augmented matrix:
  \[
    \left[\begin{array}{ccc|ccc}
      2 & 1 & 1 & 1 & 0 & 0 \\
      0 & 1 & 2 & 0 & 1 & 0 \\
      0 & 0 & 2 & 0 & 0 & 1
    \end{array}\right]
    \xrightarrow{R_1 \leftarrow \frac{1}{2} R_1,\; R_3 \leftarrow \frac{1}{2} R_3}
    \left[\begin{array}{ccc|ccc}
      1 & 1/2 & 1/2 & 1/2 & 0 & 0 \\
      0 & 1 & 2 & 0 & 1 & 0 \\
      0 & 0 & 1 & 0 & 0 & 1/2
    \end{array}\right]
  \]
  Continue to reduced row echelon form:
  \[
    \longrightarrow
    \left[\begin{array}{ccc|ccc}
      1 & 0 & 0 & 1/2 & -1/2 & 3/4 \\
      0 & 1 & 0 & 0 & 1 & -1 \\
      0 & 0 & 1 & 0 & 0 & 1/2
    \end{array}\right]
  \]
  The right block is
  \[
    \mathbf{A}^{-1} = \begin{bmatrix} 1/2 & -1/2 & 3/4 \\ 0 & 1 & -1 \\ 0 & 0 & 1/2 \end{bmatrix}.
  \]
\end{example}

\subsection{The Invertible Matrix Theorem}

\begin{theorem}{Invertible Matrix Theorem}{imt}
Let $\mathbf{A} \in \R^{n \times n}$. The following statements are equivalent:
\begin{enumerate}[label=(\alph*)]
  \item $\mathbf{A}$ is invertible.
  \item $\mathbf{A}$ is row equivalent to $\mathbf{I}_n$ (i.e.\ can be row reduced to $\mathbf{I}_n$).
  \item $\mathbf{A}$ has $n$ pivots.
  \item $\mathbf{A}\mathbf{x} = \mathbf{0}$ has only the trivial solution.
  \item The columns of $\mathbf{A}$ are linearly independent.
  \item The linear transformation $T(\mathbf{x}) = \mathbf{A}\mathbf{x}$ is one-to-one.
  \item $\mathbf{A}\mathbf{x} = \mathbf{b}$ has at least one solution for every $\mathbf{b} \in \R^n$.
  \item $\operatorname{span}\{\text{columns of }\mathbf{A}\} = \R^n$.
  \item The linear transformation $T(\mathbf{x}) = \mathbf{A}\mathbf{x}$ is onto $\R^n$.
  \item There exists $\mathbf{C}$ such that $\mathbf{C}\mathbf{A} = \mathbf{I}_n$.
  \item There exists $\mathbf{B}$ such that $\mathbf{A}\mathbf{B} = \mathbf{I}_n$.
  \item $\mathbf{A}^T$ is invertible.
\end{enumerate}
\end{theorem}

\begin{remark}
A proof sketch establishes the cycle
\[
  (a) \Longrightarrow (d) \Longrightarrow (c) \Longrightarrow (b) \Longrightarrow (a),
\]
which makes (a)--(d) equivalent. The remaining statements connect to these by similar implications.
\end{remark}

\subsection{Linear Subspace of $\R^n$}

\begin{definition}{Linear Subspace}{lin-subspace}
A subset $V \subseteq \R^n$ is a \textbf{linear subspace} if:
\begin{enumerate}
  \item $\mathbf{0} \in V$ (contains the zero vector);
  \item $\forall \mathbf{v}_1, \mathbf{v}_2 \in V$, $\mathbf{v}_1 + \mathbf{v}_2 \in V$ (closed under addition);
  \item $\forall c \in \R$, $\mathbf{v} \in V$, $c\mathbf{v} \in V$ (closed under scalar multiplication).
\end{enumerate}
\end{definition}

\begin{remark}
Condition 3 implies condition 1 (take $c = 0$) provided $V$ is non-empty; condition 1 is included for emphasis.
\end{remark}

\begin{example}
  Let
  \[
    V = \left\{ \begin{bmatrix} x_1 \\ x_2 \\ x_3 \end{bmatrix} \;\Bigm|\; x_1 + x_2 - 2x_3 = 1 \right\} \subseteq \R^3.
  \]
  Then $\mathbf{0} \notin V$, so $V$ is \textbf{not} a linear subspace. ($V$ is a plane that does not pass through the origin.)
\end{example}

\begin{example}
  Let
  \[
    V = \left\{ \begin{bmatrix} x_1 \\ x_2 \\ x_3 \end{bmatrix} \;\Bigm|\; x_1 + x_2 - 2x_3 = 0 \right\} \subseteq \R^3.
  \]
  Take $\mathbf{a} = \begin{bmatrix} a_1 \\ a_2 \\ a_3 \end{bmatrix}$, $\mathbf{b} = \begin{bmatrix} b_1 \\ b_2 \\ b_3 \end{bmatrix} \in V$.

  \textbf{Closed under addition:}
  \[
    (a_1 + b_1) + (a_2 + b_2) - 2(a_3 + b_3) = (a_1 + a_2 - 2a_3) + (b_1 + b_2 - 2b_3) = 0,
  \]
  so $\mathbf{a} + \mathbf{b} = \begin{bmatrix} a_1 + b_1 \\ a_2 + b_2 \\ a_3 + b_3 \end{bmatrix} \in V$.

  \textbf{Closed under scalar multiplication:} for $c \in \R$,
  \[
    c x_1 + c x_2 - 2 c x_3 = c(x_1 + x_2 - 2 x_3) = c \cdot 0 = 0,
  \]
  so $c\mathbf{a} \in V$. Combined with $\mathbf{0} \in V$, $V$ is a linear subspace.
\end{example}

\subsubsection{Span is Linear Subspace}

For $\mathbf{v}_1, \mathbf{v}_2 \in \R^n$,
\[
  V = \operatorname{span}\{\mathbf{v}_1, \mathbf{v}_2\} = \{\, c_1 \mathbf{v}_1 + c_2 \mathbf{v}_2 \mid c_1, c_2 \in \R \,\}.
\]

\begin{theorem}{Span Is a Linear Subspace}{span-subspace}
For all $\mathbf{v}_1, \mathbf{v}_2 \in \R^n$, $V = \operatorname{span}\{\mathbf{v}_1, \mathbf{v}_2\}$ is a linear subspace of $\R^n$.
\end{theorem}

\begin{proof}
\textbf{(1) Closed under addition.} Let $\mathbf{w}_1 = c_1 \mathbf{v}_1 + c_2 \mathbf{v}_2$, $\mathbf{w}_2 = d_1 \mathbf{v}_1 + d_2 \mathbf{v}_2 \in V$. Then
\[
  \mathbf{w}_1 + \mathbf{w}_2
  = (c_1 \mathbf{v}_1 + c_2 \mathbf{v}_2) + (d_1 \mathbf{v}_1 + d_2 \mathbf{v}_2)
  = (c_1 + d_1) \mathbf{v}_1 + (c_2 + d_2) \mathbf{v}_2
\]
is a linear combination of $\mathbf{v}_1, \mathbf{v}_2$, so $\mathbf{w}_1 + \mathbf{w}_2 \in V$.

\textbf{(2) Closed under scalar multiplication.} Let $\mathbf{w} = c_1 \mathbf{v}_1 + c_2 \mathbf{v}_2 \in V$. For $c \in \R$,
\[
  c\mathbf{w} = c(c_1 \mathbf{v}_1 + c_2 \mathbf{v}_2) = (c c_1) \mathbf{v}_1 + (c c_2) \mathbf{v}_2
\]
is a linear combination of $\mathbf{v}_1, \mathbf{v}_2$, so $c\mathbf{w} \in V$.
\end{proof}

\begin{theorem}{Generalized Span Is a Linear Subspace}{span-subspace-gen}
For $\mathbf{v}_1, \mathbf{v}_2, \ldots, \mathbf{v}_k \in \R^n$, $\operatorname{span}\{\mathbf{v}_1, \ldots, \mathbf{v}_k\}$ is a linear subspace of $\R^n$.
\end{theorem}

\subsubsection{Null Space and Column Space}

Let $\mathbf{A} \in \R^{m \times n}$.

\begin{definition}{Null Space}{null-space}
The \textbf{null space} of $\mathbf{A}$ is
\[
  \operatorname{Null}(\mathbf{A}) = \{\, \mathbf{x} \in \R^n \mid \mathbf{A}\mathbf{x} = \mathbf{0} \,\}.
\]
(The solution set of the homogeneous system $\mathbf{A}\mathbf{x} = \mathbf{0}$.)
\end{definition}

\begin{definition}{Column Space}{col-space}
Writing $\mathbf{A} = [\,\mathbf{a}_1 \mid \mathbf{a}_2 \mid \cdots \mid \mathbf{a}_n\,]$, the \textbf{column space} of $\mathbf{A}$ is
\[
  \operatorname{Col}(\mathbf{A}) = \operatorname{span}\{\mathbf{a}_1, \mathbf{a}_2, \ldots, \mathbf{a}_n\}.
\]
\end{definition}

\begin{theorem}{Null and Column Spaces Are Subspaces}{null-col-subspace}
$\operatorname{Null}(\mathbf{A})$ is a linear subspace of $\R^n$, and $\operatorname{Col}(\mathbf{A})$ is a linear subspace of $\R^m$.
\end{theorem}

\begin{proof}
$\operatorname{Col}(\mathbf{A})$ is a span, hence a linear subspace by Theorem~\ref{thm:span-subspace-gen}.

For $\operatorname{Null}(\mathbf{A})$:

\textbf{(1) Closed under addition.} For $\mathbf{x}_1, \mathbf{x}_2 \in \operatorname{Null}(\mathbf{A})$, $\mathbf{A}\mathbf{x}_1 = \mathbf{0}$ and $\mathbf{A}\mathbf{x}_2 = \mathbf{0}$, so
\[
  \mathbf{A}(\mathbf{x}_1 + \mathbf{x}_2) = \mathbf{A}\mathbf{x}_1 + \mathbf{A}\mathbf{x}_2 = \mathbf{0},
\]
hence $\mathbf{x}_1 + \mathbf{x}_2 \in \operatorname{Null}(\mathbf{A})$.

\textbf{(2) Closed under scalar multiplication.} For $c \in \R$ and $\mathbf{x} \in \operatorname{Null}(\mathbf{A})$,
\[
  \mathbf{A}(c\mathbf{x}) = c(\mathbf{A}\mathbf{x}) = c \cdot \mathbf{0} = \mathbf{0},
\]
hence $c\mathbf{x} \in \operatorname{Null}(\mathbf{A})$.
\end{proof}

\subsection{Basis of a Linear Subspace}

Let $V \subseteq \R^n$ be a linear subspace.

\begin{definition}{Basis}{basis}
A set $\{\mathbf{v}_1, \mathbf{v}_2, \ldots, \mathbf{v}_k\} \subseteq V$ is a \textbf{basis} of $V$ if:
\begin{enumerate}
  \item $\operatorname{span}\{\mathbf{v}_1, \ldots, \mathbf{v}_k\} = V$;
  \item $\{\mathbf{v}_1, \mathbf{v}_2, \ldots, \mathbf{v}_k\}$ are linearly independent.
\end{enumerate}
\end{definition}

\begin{example}
  Let $\mathbf{A} = \begin{bmatrix} 1 & 2 & 3 \\ 2 & 4 & 6 \end{bmatrix}$.
  \[
    \operatorname{Col}(\mathbf{A})
    = \operatorname{span}\left\{ \begin{bmatrix} 1 \\ 2 \end{bmatrix}, \begin{bmatrix} 2 \\ 4 \end{bmatrix}, \begin{bmatrix} 3 \\ 6 \end{bmatrix} \right\}
    = \operatorname{span}\left\{ \begin{bmatrix} 1 \\ 2 \end{bmatrix} \right\},
  \]
  so $\left\{ \begin{bmatrix} 1 \\ 2 \end{bmatrix} \right\}$ is a basis for $\operatorname{Col}(\mathbf{A})$.

  For the null space, $\mathbf{A}\mathbf{x} = \mathbf{0}$ reduces to $x_1 + 2 x_2 + 3 x_3 = 0$, so $x_1 = -2 x_2 - 3 x_3$ with $x_2, x_3$ free. Then
  \[
    \operatorname{Null}(\mathbf{A})
    = \operatorname{span}\left\{ \begin{bmatrix} -2 \\ 1 \\ 0 \end{bmatrix}, \begin{bmatrix} -3 \\ 0 \\ 1 \end{bmatrix} \right\},
  \]
  and these two vectors form a basis for $\operatorname{Null}(\mathbf{A})$.
\end{example}

\subsection{Dimension}

\begin{theorem}{Bases Have the Same Size}{basis-size}
If $B$ and $B'$ are both bases of a linear subspace $V$, then $B$ and $B'$ have the same number of vectors.
\end{theorem}

\begin{definition}{Dimension}{dimension}
Let $V$ be a linear subspace of $\R^n$ with basis $B$. The \textbf{dimension} of $V$, written $\dim(V)$, is the number of vectors in $B$.
\end{definition}

\begin{example}
  For $\mathbf{A} = \begin{bmatrix} 1 & 2 & 3 \\ 2 & 4 & 6 \end{bmatrix}$, $\dim(\operatorname{Null}(\mathbf{A})) = 2$.
\end{example}

\begin{example}
  $\dim(\R^n) = n$, with standard basis $\mathbf{e}_1, \mathbf{e}_2, \ldots, \mathbf{e}_n$.
\end{example}

\subsection{Rank-Nullity Theorem}

\begin{theorem}{Rank-Nullity}{rank-nullity}
Let $\mathbf{A} \in \R^{m \times n}$. Then
\[
  \dim(\operatorname{Null}(\mathbf{A})) + \dim(\operatorname{Col}(\mathbf{A})) = n.
\]
\end{theorem}

\begin{definition}{Rank and Nullity}{rank-nullity-def}
For $\mathbf{A} \in \R^{m \times n}$,
\[
  \operatorname{nullity}(\mathbf{A}) = \dim(\operatorname{Null}(\mathbf{A})),
  \qquad
  \operatorname{rank}(\mathbf{A}) = \dim(\operatorname{Col}(\mathbf{A})).
\]
\end{definition}

\begin{example}
  Let $\mathbf{A} = \begin{bmatrix} 1 & 2 & 3 \\ 2 & 4 & 6 \end{bmatrix} \in \R^{2 \times 3}$. From the previous bases, $\operatorname{rank}(\mathbf{A}) = 1$ and $\operatorname{nullity}(\mathbf{A}) = 2$, and $1 + 2 = 3 = n$.
\end{example}

\begin{example}
  Let
  \[
    \mathbf{A} = \begin{bmatrix} 1 & 2 & -3 & 4 & 1 \\ 2 & -1 & 3 & 1 & 0 \\ 3 & 1 & 0 & 5 & 1 \end{bmatrix} \in \R^{3 \times 5}.
  \]
  Row reduce:
  \[
    \mathbf{A}
    \;\rightsquigarrow\;
    \begin{bmatrix} 1 & 2 & -3 & 4 & 1 \\ 0 & -5 & 9 & -7 & -2 \\ 0 & 0 & 0 & 0 & 0 \end{bmatrix}.
  \]
  Pivots are in columns $1$ and $2$, so $\operatorname{rank}(\mathbf{A}) = 2$ and by Theorem~\ref{thm:rank-nullity}, $\operatorname{nullity}(\mathbf{A}) = 5 - 2 = 3$.
\end{example}
